\section{Introduction}
\label{sec:introduction}

\subsection{The Large-$k$ Mixing Problem}

Standard \recom{}~\citep{deford2021} mixes well for small district counts
($k \leq 14$, e.g.\ North Carolina), but degrades for the large states that
matter most for congressional representation.
Texas ($k = 38$) and California ($k = 52$) each require an ensemble that
explores meaningfully different plans---plans that differ in which counties
anchor which districts, not just which tracts lie on which boundaries.
A \recom{} step merges two randomly chosen adjacent districts and redraws
the boundary via a random spanning tree cut.
With $k = 38$ districts, a single step touches $2/38 \approx 5\%$ of the
plan.
To reach a plan that differs substantially from the start---one where, say,
the Dallas metro is split differently from the Harris County anchor---the
chain must restructure large contiguous areas, which requires $O(k)$ consecutive
steps in the same geographic region.
In practice, auto-correlation in the edge-cut statistic decays slowly: at
lag 100, standard \recom{} on Texas still exhibits correlations above 0.5,
indicating that 100-step windows are far too short to treat as independent
draws from the plan space~\citep{autry2021}.

This slow mixing is not a convergence failure---\recom{} does converge, and
given sufficient steps it samples the plan space correctly---but a practical
obstacle.
Ensemble analysis of redistricting plans for litigation, academic study, or
legislative compliance requires samples that are effectively independent.
For TX and CA, achieving independence requires impractically long chains
with single-scale \recom{}.

\subsection{Multi-scale MCMC as the Hierarchical Solution}

Multi-scale MCMC~\citep{autry2021} resolves the large-$k$ mixing problem
by operating simultaneously at two spatial resolutions.
A \emph{coarse} chain restructures the plan at county granularity---treating
each county as an atomic unit and proposing ReCom moves on the county
adjacency graph.
A \emph{fine} chain refines boundaries at tract granularity, moving
individual tracts between adjacent districts.
At each step, the algorithm selects the coarse chain with probability
$\alpha$ (the coarsening parameter) and the fine chain with probability
$1 - \alpha$.

Coarse moves are transformative: a single ReCom step on the county graph
can simultaneously reassign dozens of tracts across a large geographic area,
achieving in one step what would take $O(n_{\mathrm{county}})$ fine steps.
Fine moves are precise: they refine district boundaries at the tract level,
achieving good population balance and compact boundaries within the
constraint set by the coarse structure.
The interleaving of the two scales achieves faster global mixing (via coarse
moves) without sacrificing local boundary precision (via fine moves).

\subsection{Three Contributions}

This paper makes three contributions.

\begin{enumerate}

  \item \textbf{County-level coarsening without extra data.}
  We show that the county adjacency graph required for coarse \recom{} moves
  can be derived exactly from the existing tract adjacency graph using GEOID
  prefix matching.
  No extra data files, shapefiles, or census downloads are required.
  Section~\ref{sec:algorithm} gives the derivation and proves its exactness
  (Proposition~\ref{prop:county-adj-exact}).

  \item \textbf{Implementation of Option B.}
  We implement \ms{} Option B (fine = census tract, coarse = county) in
  \texttt{redist-ensemble} as \texttt{MultiScaleChain} and expose it via
  \texttt{--search multiscale} in the \texttt{redist-cli}.
  Section~\ref{sec:algorithm} gives the complete algorithm specification,
  including the rebalance contract (200 boundary swaps, reject if fails)
  and the coarse population tolerance ($3\times$ the fine tolerance).

  \item \textbf{Mixing speed comparison.}
  We compare autocorrelation at lag 100 between standard \recom{} and
  \ms{} (Option B, $\alpha = 0.3$) on Texas ($k = 38$, $T = 2000$ steps).
  Section~\ref{sec:mixing} reports the results; the expected reduction is
  $40$--$50\%$ in lag-100 autocorrelation at comparable wall-clock time.

\end{enumerate}

\subsection{Relation to Prior Work}

Autry et al.~\citep{autry2021} introduced multi-scale MCMC for redistricting
and proved that the combined chain is ergodic under mild conditions.
Their empirical results on North Carolina ($k = 13$) showed a $3\times$
speedup in effective sample rate relative to single-scale \recom{}.
Our contribution is an implementation that derives the coarse graph from
existing data structures (no extra preprocessing) and validates the approach
on the larger states (TX, CA) where mixing improvement is most needed.

DeFord et al.~\citep{deford2021} introduced the \recom{} proposal that forms
the fine-level chain in our implementation.
The coarse-level chain uses the same ReCom mechanism on the county graph,
so both scales share the same sampling subroutine.

The resolution system (Section~\ref{sec:resolution}) that unifies Options A,
B, and C is described in the companion resolution spec~\citep{dellaLibera2026resolution}.
Forest ReCom~\citep{dellaLibera2026forestRecom} (G.9) provides an alternative
fine-level chain with a known stationary distribution; it can replace standard
ReCom in the fine level of \ms{} for distributional correctness, though this
extension is left for future work.

\subsection{Paper Structure}

Section~\ref{sec:background} reviews standard \recom{} mixing and the
theory of coarse-level abstraction.
Section~\ref{sec:algorithm} gives the formal specification of \ms{} Option B.
Section~\ref{sec:resolution} describes the three resolution options (A, B, C)
and the shared \texttt{derive\_partition} helper.
Section~\ref{sec:mixing} reports mixing speed results on Texas.
Section~\ref{sec:conclusion} concludes with guidelines for when to use \ms{}.
